\chapter{Enums}\label{ch:enums}
\index{enum}

A struct says ``this value has \emph{all} of these fields.'' An enum says
``this value is \emph{one} of these variants.'' Where structs model
\emph{product types}---a user has a name \emph{and} an age \emph{and} an
email---enums model \emph{sum types}: a network response is either a success
\emph{or} a failure, a shape is a circle \emph{or} a rectangle \emph{or}
a triangle.

Lattice enums are full algebraic data types. Each variant can carry its own
payload, and the compiler checks that your \lstinline|match| expressions
handle every variant. In this chapter we will define enums, construct
variants, destructure them with pattern matching, model state machines,
and build the ubiquitous \lstinline|Option| and \lstinline|Result| patterns.

%% ─────────────────────────────────────────────────────────────────────
\section{Defining Enums With and Without Payloads}\label{sec:enum-defining}
\index{enum!definition}
\index{enum!variant}

An enum definition lists a type name and a set of \emph{variants}:

\begin{lstlisting}[language=Lattice, caption={A basic enum with unit variants}]
enum Direction {
    North,
    South,
    East,
    West
}
\end{lstlisting}

Each variant is a distinct tag. A \lstinline|Direction| value is exactly
one of \lstinline|North|, \lstinline|South|, \lstinline|East|, or
\lstinline|West|---nothing else.

\begin{defnbox}{Enum}
An \emph{enum} (short for ``enumeration'') is a type that can be one of
a fixed set of \emph{variants}. Each variant has a name and may carry a
\emph{payload}---zero or more associated values that travel with the tag.
\end{defnbox}

Variants without payloads are called \emph{unit variants}. They carry no
additional data:

\begin{lstlisting}[language=Lattice, caption={Unit variants carry no data}]
enum Color {
    Red,
    Green,
    Blue
}

let favourite = Color::Red
print(favourite)  // Color::Red
\end{lstlisting}

\subsection{Variants with Payloads}

A variant becomes more interesting when it carries data. List the payload
types in parentheses after the variant name:

\begin{lstlisting}[language=Lattice, caption={Enum with payload variants}]
enum Color {
    Red,
    Green,
    Blue,
    Rgb(any, any, any)
}

let coral = Color::Rgb(255, 127, 80)
print(coral)  // Color::Rgb(255, 127, 80)
\end{lstlisting}

The \lstinline|Rgb| variant carries three values---the red, green, and blue
channels. Notice that unit variants and payload variants can coexist in the
same enum. You can mix and match freely.

Here is a \lstinline|Shape| enum where each variant carries different amounts
of data:

\begin{lstlisting}[language=Lattice, caption={Variants with different payload sizes}]
enum Shape {
    Circle(any),
    Rectangle(any, any),
    Triangle(any, any, any),
    None
}
\end{lstlisting}

\lstinline|Circle| takes one value (the radius), \lstinline|Rectangle| takes
two (width and height), \lstinline|Triangle| takes three (side lengths),
and \lstinline|None| carries nothing.

\begin{notebox}{Payload Type Annotations}
In the examples above we use \lstinline|any| for payload types. You can also
write concrete types for documentation and strict-mode checking:

\begin{lstlisting}[language=Lattice]
enum Shape {
    Circle(Float),
    Rectangle(Float, Float),
    None
}
\end{lstlisting}

This makes the intended contract clearer. Generic type parameters (e.g.,
\lstinline|Option<T>|) are covered in \cref{ch:generics}.
\end{notebox}

%% ─────────────────────────────────────────────────────────────────────
\section{Constructing Variants}\label{sec:enum-construction}
\index{enum!construction}

To create an enum value, use the double-colon syntax:
\lstinline|EnumName::VariantName| for unit variants, or
\lstinline|EnumName::VariantName(args)| for payload variants.

\begin{lstlisting}[language=Lattice, caption={Constructing enum variants}]
enum Color {
    Red,
    Green,
    Blue,
    Rgb(any, any, any)
}

let r = Color::Red
let g = Color::Green
let custom = Color::Rgb(255, 128, 0)

print(r)       // Color::Red
print(g)       // Color::Green
print(custom)  // Color::Rgb(255, 128, 0)
\end{lstlisting}

\subsection{Enums Returned from Functions}

Enums are regular values---you can pass them to functions, return them, and
store them in collections:

\begin{lstlisting}[language=Lattice, caption={Returning enums from functions}]
enum Color {
    Red,
    Green,
    Blue,
    Rgb(any, any, any)
}

fn make_color(kind: String) -> Color {
    if kind == "red" {
        return Color::Red
    }
    if kind == "custom" {
        return Color::Rgb(1, 2, 3)
    }
    return Color::Blue
}

print(make_color("red"))     // Color::Red
print(make_color("custom"))  // Color::Rgb(1, 2, 3)
print(make_color("other"))   // Color::Blue
\end{lstlisting}

\subsection{Enums in Collections}

\begin{lstlisting}[language=Lattice, caption={An array of enum values}]
enum Status {
    Pending,
    Running(any),
    Done(any),
    Failed(any, any)
}

let tasks = [
    Status::Pending,
    Status::Running("compile"),
    Status::Done(42),
    Status::Failed("timeout", 504)
]

for task in tasks {
    print(task)
}
\end{lstlisting}

\subsection{Runtime Introspection Methods}
\index{enum!tag}
\index{enum!payload}

Every enum value carries metadata that you can inspect at runtime:

\begin{lstlisting}[language=Lattice, caption={Enum introspection methods}]
enum Option {
    Some(any),
    None
}

let val = Option::Some(42)

print(val.tag())           // "Some"
print(val.enum_name())     // "Option"
print(val.variant_name())  // "Some"
print(val.payload())       // [42]
print(val.is_variant("Some"))  // true
print(val.is_variant("None"))  // false
\end{lstlisting}

The \lstinline|.payload()| method returns an array of the variant's payload
elements. For unit variants, it returns an empty array. These methods are
useful for logging and debugging, but for most logic you should prefer
pattern matching.

%% ─────────────────────────────────────────────────────────────────────
\section{Matching on Enums}\label{sec:enum-matching}
\index{enum!matching}
\index{match!enum patterns}
\index{pattern matching!enum}

Pattern matching is the primary way to work with enums. If you read
\cref{ch:pattern-matching}, you have already seen \lstinline|match|
expressions with literals, wildcards, and guards. Enum patterns add a new
dimension: you can match on the variant \emph{and} destructure the payload
in a single step.

\subsection{Matching Unit Variants}

At its most basic, matching an enum looks like matching constants:

\begin{lstlisting}[language=Lattice, caption={Matching unit variants}]
enum Direction {
    North,
    South,
    East,
    West
}

fn describe(dir: Direction) -> String {
    match dir {
        Direction::North => "heading north",
        Direction::South => "heading south",
        Direction::East  => "heading east",
        Direction::West  => "heading west"
    }
}

print(describe(Direction::North))  // heading north
print(describe(Direction::West))   // heading west
\end{lstlisting}

\subsection{Destructuring Payloads}

When a variant carries data, the match pattern can bind payload elements
to variables:

\begin{lstlisting}[language=Lattice, caption={Destructuring enum payloads}]
enum Option {
    Some(any),
    None
}

let val = Option::Some(42)

let result = match val {
    Option::Some(x) => "got: " + to_string(x),
    Option::None    => "nothing",
    _ => "unknown"
}
print(result)  // got: 42
\end{lstlisting}

The pattern \lstinline|Option::Some(x)| does two things at once: it checks
that the value is the \lstinline|Some| variant, and it binds the payload to
the variable \lstinline|x|. Inside the arm's body, \lstinline|x| holds the
value \lstinline|42|.

\subsection{Multi-Element Payloads}

Variants with multiple payload elements bind multiple variables:

\begin{lstlisting}[language=Lattice, caption={Multi-binding destructuring}]
enum Pair {
    Both(any, any),
    Single(any),
    Empty
}

let p = Pair::Both("hello", 99)
let result = match p {
    Pair::Both(a, b)  => to_string(a) + "=" + to_string(b),
    Pair::Single(x)   => "single: " + to_string(x),
    Pair::Empty        => "empty",
    _ => "?"
}
print(result)  // hello=99
\end{lstlisting}

\subsection{Using Destructured Values in Computations}

The bound variables are full values---you can use them in any expression:

\begin{lstlisting}[language=Lattice, caption={Computing with destructured values}]
enum Result {
    Ok(any),
    Err(any)
}

let r = Result::Ok(21)
let doubled = match r {
    Result::Ok(v)  => v * 2,
    Result::Err(e) => -1,
    _ => 0
}
print(doubled)  // 42
\end{lstlisting}

\subsection{Block Bodies in Arms}

When a single expression is not enough, use a block body:

\begin{lstlisting}[language=Lattice, caption={Block bodies in match arms}]
enum Option {
    Some(any),
    None
}

let val = Option::Some(10)
let result = match val {
    Option::Some(n) => {
        let doubled = n * 2
        let tripled = n * 3
        "doubled=" + to_string(doubled) + " tripled=" + to_string(tripled)
    },
    _ => "none"
}
print(result)  // doubled=20 tripled=30
\end{lstlisting}

\subsection{Ignoring Payload with Wildcards}

If you care about the variant but not the payload, use \lstinline|_| in the
payload position:

\begin{lstlisting}[language=Lattice, caption={Ignoring payloads with wildcards}]
enum Status {
    Pending,
    Running(any),
    Done(any),
    Failed(any, any)
}

flux pending_count = 0
flux done_count = 0
let statuses = [Status::Pending, Status::Done(42), Status::Running("build")]

for s in statuses {
    match s {
        Status::Pending => { pending_count += 1 },
        Status::Done(_) => { done_count += 1 },
        _ => {}
    }
}
print(pending_count)  // 1
print(done_count)     // 1
\end{lstlisting}

\subsection{Nested Match Expressions}

You can nest match expressions to drill into layered data:

\begin{lstlisting}[language=Lattice, caption={Nested match on enums}]
enum Result {
    Ok(any),
    Err(any)
}

fn classify(val: any) -> String {
    return match val {
        Result::Ok(v) => {
            match v {
                0 => "ok:zero",
                _ => "ok:" + to_string(v)
            }
        },
        Result::Err(e) => "err:" + to_string(e),
        _ => "unknown"
    }
}

print(classify(Result::Ok(0)))       // ok:zero
print(classify(Result::Ok(42)))      // ok:42
print(classify(Result::Err("bad")))  // err:bad
\end{lstlisting}

\begin{tipbox}{Match as an Expression}
Remember that \lstinline|match| is an expression in Lattice---it evaluates
to a value. You can assign its result to a variable, pass it to a function,
or return it directly. This makes match-based dispatch composable and
concise.
\end{tipbox}

%% ─────────────────────────────────────────────────────────────────────
\section{Using Enums for State Machines}\label{sec:enum-state-machines}
\index{enum!state machine}
\index{state machine}

Enums are a natural fit for state machines because each variant represents
a distinct state, and the payload carries state-specific data. Let's model
a task runner that moves through a lifecycle:

\begin{lstlisting}[language=Lattice, caption={A task lifecycle as an enum}]
enum TaskState {
    Pending,
    Running(any),
    Done(any),
    Failed(any, any)
}

fn status_label(s: TaskState) -> String {
    match s {
        TaskState::Pending => "pending",
        TaskState::Running(task) =>
            "running: " + to_string(task),
        TaskState::Done(result) =>
            "done: " + to_string(result),
        TaskState::Failed(err, code) =>
            "failed: " + to_string(err) + " (" + to_string(code) + ")",
        _ => "unknown"
    }
}
\end{lstlisting}

\subsection{Driving State Transitions}

A state machine advances by replacing the current state with a new variant:

\begin{lstlisting}[language=Lattice, caption={State transitions}]
enum TaskState {
    Pending,
    Running(any),
    Done(any),
    Failed(any, any)
}

fn advance(state: TaskState, tick: Int) -> TaskState {
    match state {
        TaskState::Pending => TaskState::Running("job_" + to_string(tick)),
        TaskState::Running(task) => {
            if tick > 3 {
                TaskState::Done("completed " + to_string(task))
            } else {
                TaskState::Running(task)
            }
        },
        _ => state
    }
}

flux state = TaskState::Pending
for tick in 0..6 {
    state = advance(state, tick)
    print(to_string(tick) + ": " + match state {
        TaskState::Pending      => "pending",
        TaskState::Running(t)   => "running " + to_string(t),
        TaskState::Done(r)      => "done: " + to_string(r),
        TaskState::Failed(e, c) => "failed",
        _ => "?"
    })
}
\end{lstlisting}

Because the state is an enum, the compiler can warn you if you forget to
handle a variant. And because each variant carries exactly the data it needs,
you never have stale or irrelevant fields cluttering your state object.

\subsection{A Parser Token Example}

Enums excel at modelling token types in a parser, where each token variant
carries different data:

\begin{lstlisting}[language=Lattice, caption={Tokens as an enum}]
enum Token {
    Number(any),
    Identifier(any),
    Plus,
    Minus,
    Star,
    LeftParen,
    RightParen,
    Eof
}

fn token_label(tok: Token) -> String {
    match tok {
        Token::Number(n)     => "NUM(" + to_string(n) + ")",
        Token::Identifier(s) => "ID(" + to_string(s) + ")",
        Token::Plus          => "+",
        Token::Minus         => "-",
        Token::Star          => "*",
        Token::LeftParen     => "(",
        Token::RightParen    => ")",
        Token::Eof           => "EOF"
    }
}

let tokens = [
    Token::Number(42),
    Token::Plus,
    Token::Identifier("x"),
    Token::Eof
]

for tok in tokens {
    print(token_label(tok))
}
// Output: NUM(42), +, ID(x), EOF
\end{lstlisting}

%% ─────────────────────────────────────────────────────────────────────
\section{Option-Like and Result-Like Patterns}\label{sec:enum-option-result}
\index{Option}
\index{Result}
\index{enum!Option pattern}
\index{enum!Result pattern}

Two patterns appear so frequently across programming languages that they
deserve special attention: \emph{optional values} (a value that might be
absent) and \emph{fallible operations} (a computation that might fail).
Lattice does not have built-in \lstinline|Option| and \lstinline|Result|
types in the language itself, but you can define them as enums and they
work beautifully.

\subsection{The Option Pattern}

\begin{lstlisting}[language=Lattice, caption={Defining and using Option}]
enum Option {
    Some(any),
    None
}

fn find_user(id: Int) -> Option {
    if id == 1 {
        return Option::Some("Alice")
    }
    if id == 2 {
        return Option::Some("Bob")
    }
    return Option::None
}

let user = find_user(1)
match user {
    Option::Some(name) => print("Found: " + name),
    Option::None       => print("User not found"),
    _ => print("?")
}
// Output: Found: Alice
\end{lstlisting}

An \lstinline|unwrap_or| helper makes optionals easier to use in
expressions:

\begin{lstlisting}[language=Lattice, caption={An unwrap\_or helper}]
enum Option {
    Some(any),
    None
}

fn unwrap_or(opt: Option, default: any) -> any {
    match opt {
        Option::Some(v) => v,
        Option::None    => default,
        _ => default
    }
}

let a = Option::Some(10)
let b = Option::None

print(unwrap_or(a, 0))   // 10
print(unwrap_or(b, -1))  // -1
\end{lstlisting}

\subsection{The Result Pattern}

\begin{lstlisting}[language=Lattice, caption={Defining and using Result}]
enum Result {
    Ok(any),
    Err(any)
}

fn parse_port(input: String) -> Result {
    let n = parse_int(input)
    if n == nil {
        return Result::Err("not a number: " + input)
    }
    if n < 1 || n > 65535 {
        return Result::Err("port out of range: " + to_string(n))
    }
    return Result::Ok(n)
}

let result = parse_port("8080")
match result {
    Result::Ok(port)  => print("Listening on port " + to_string(port)),
    Result::Err(msg)  => print("Error: " + msg),
    _ => print("?")
}
// Output: Listening on port 8080
\end{lstlisting}

\subsection{Chaining Results}

A common pattern is to chain fallible operations---if the first succeeds,
feed its value into the second:

\begin{lstlisting}[language=Lattice, caption={Chaining two results}]
enum Result {
    Ok(any),
    Err(any)
}

fn chain_results(a: Result, b: Result) -> Result {
    match a {
        Result::Ok(va) => {
            match b {
                Result::Ok(vb) => Result::Ok(va + vb),
                Result::Err(e) => Result::Err(e),
                _ => Result::Err("unknown")
            }
        },
        Result::Err(e) => Result::Err(e),
        _ => Result::Err("unknown")
    }
}

let r = chain_results(Result::Ok(10), Result::Ok(20))
match r {
    Result::Ok(v)  => print("ok: " + to_string(v)),
    Result::Err(e) => print("err: " + to_string(e)),
    _ => print("unknown")
}
// Output: ok: 30
\end{lstlisting}

\begin{notebox}{Enums vs.\ \texttt{try}/\texttt{catch}}
Lattice also provides \lstinline|try|/\lstinline|catch| for error handling
(see \cref{ch:error-handling}). The \lstinline|Result| pattern and
exceptions are complementary strategies. Use \lstinline|Result| when errors
are \emph{expected} and part of your data flow; use
\lstinline|try|/\lstinline|catch| when errors are \emph{exceptional} and you
want to bail out of a computation.
\end{notebox}

%% ─────────────────────────────────────────────────────────────────────
\section{Enum Exhaustiveness Guarantees}\label{sec:enum-exhaustiveness}
\index{enum!exhaustiveness}
\index{match!exhaustiveness}

One of the most powerful features of enum-based matching is
\emph{exhaustiveness checking}. When you match on an enum, the compiler
examines your arms and warns you if any variant is unhandled.

\subsection{How It Works}

The compiler (in \filename{src/match\_check.c}) performs a straightforward
analysis:

\begin{enumerate}
  \item If any arm has an unguarded wildcard (\lstinline|_|) or binding
    pattern, the match is considered exhaustive---the catch-all handles
    everything.
  \item Otherwise, if the patterns are enum variants, the compiler collects
    which variants appear and compares them against the enum's definition.
  \item Any missing variants produce a warning.
\end{enumerate}

\begin{lstlisting}[language=Lattice, caption={An exhaustive match---no warning}]
enum Direction {
    North,
    South,
    East,
    West
}

fn label(d: Direction) -> String {
    match d {
        Direction::North => "N",
        Direction::South => "S",
        Direction::East  => "E",
        Direction::West  => "W"
    }
}
\end{lstlisting}

All four variants are covered, so the compiler is happy. Now remove one:

\begin{lstlisting}[language=Lattice, caption={A non-exhaustive match---compiler warns}]
fn label(d: Direction) -> String {
    match d {
        Direction::North => "N",
        Direction::South => "S",
        Direction::East  => "E"
        // West is missing!
    }
}
// Compiler warning: non-exhaustive match: missing Direction variant `Direction::West`
\end{lstlisting}

\begin{warnbox}{Guarded Arms and Exhaustiveness}
A guarded arm (e.g., \lstinline|n if n > 0 => ...|) still counts toward
variant coverage---the compiler sees that you explicitly handled that
variant. However, a guarded \emph{wildcard} does \emph{not} count as a
catch-all, because the guard might fail at runtime. If you need a true
catch-all, add an unguarded \lstinline|_ => ...| arm.
\end{warnbox}

\subsection{Boolean Exhaustiveness}

The checker also handles boolean matches:

\begin{lstlisting}[language=Lattice, caption={Boolean exhaustiveness}]
let flag = true
let label = match flag {
    true  => "yes",
    false => "no"
}
// No warning: both cases covered
\end{lstlisting}

If you omit \lstinline|true| or \lstinline|false|, the compiler warns about
the missing case.

\subsection{Non-Enum Matches}

For integer, string, or float scrutinees, the set of possible values is
unbounded. The compiler cannot enumerate them, so it requires a wildcard arm:

\begin{lstlisting}[language=Lattice, caption={Wildcard required for integer match}]
let score = 85
let grade = match score {
    0..59   => "F",
    60..69  => "D",
    70..79  => "C",
    80..89  => "B",
    90..100 => "A",
    _       => "invalid"
}
print(grade)  // B
\end{lstlisting}

Without the \lstinline|_ => "invalid"| arm, the compiler warns:
``consider adding a wildcard \lstinline|_| arm.''

\subsection{Why Exhaustiveness Matters}

Exhaustiveness checking catches a common class of bugs at compile time.
When you add a new variant to an enum, every \lstinline|match| expression
that touches that enum will produce a warning until you handle the new case.
This is far safer than an \lstinline|if|/\lstinline|else| chain, where a
missing branch silently falls through.

\begin{tipbox}{The ``Add a Variant'' Test}
When designing an enum, imagine adding a new variant six months from now.
Will all the match sites in your codebase light up with warnings? If you
use enums and match consistently, the answer is yes---and that is a
tremendous safety net for evolving codebases.
\end{tipbox}

%% ─────────────────────────────────────────────────────────────────────
\section{Recursive Enums}\label{sec:enum-recursive}
\index{enum!recursive}

Because payload elements can be any value, enums can refer to themselves
recursively. This is perfect for tree structures:

\begin{lstlisting}[language=Lattice, caption={A recursive binary tree}]
enum Tree {
    Leaf(any),
    Branch(any, any)
}

fn tree_sum(t: any) -> Int {
    match t {
        Tree::Leaf(v) => v,
        Tree::Branch(l, r) => tree_sum(l) + tree_sum(r),
        _ => 0
    }
}

fn tree_depth(t: any) -> Int {
    match t {
        Tree::Leaf(_) => 1,
        Tree::Branch(l, r) => {
            let ld = tree_depth(l)
            let rd = tree_depth(r)
            if ld > rd { 1 + ld } else { 1 + rd }
        },
        _ => 0
    }
}

fn make_balanced(depth: Int, value: Int) -> Tree {
    if depth <= 0 {
        return Tree::Leaf(value)
    }
    return Tree::Branch(
        make_balanced(depth - 1, value * 2),
        make_balanced(depth - 1, value * 2 + 1)
    )
}

let tree = make_balanced(3, 1)
print("sum: " + to_string(tree_sum(tree)))
print("depth: " + to_string(tree_depth(tree)))
\end{lstlisting}

Each \lstinline|Branch| holds two children that are themselves
\lstinline|Tree| values. The recursive functions use match to distinguish
leaves from branches, and the wildcard arm handles unexpected cases
gracefully.

\subsection{An Expression Evaluator}

Recursive enums can model abstract syntax trees for a mini-language:

\begin{lstlisting}[language=Lattice, caption={An expression tree evaluator}]
enum Expr {
    Num(any),
    Add(any, any),
    Mul(any, any),
    Neg(any)
}

fn eval_expr(e: Expr) -> Int {
    match e {
        Expr::Num(n)    => n,
        Expr::Add(a, b) => eval_expr(a) + eval_expr(b),
        Expr::Mul(a, b) => eval_expr(a) * eval_expr(b),
        Expr::Neg(x)    => -eval_expr(x),
        _ => 0
    }
}

// Represents: (3 * 4) + (-(2))
let expr = Expr::Add(
    Expr::Mul(Expr::Num(3), Expr::Num(4)),
    Expr::Neg(Expr::Num(2))
)
print(eval_expr(expr))  // 10
\end{lstlisting}

This is a pattern you will encounter throughout compiler design, symbolic
mathematics, and rule engines. Enums make the variant structure explicit;
match makes the dispatch exhaustive.

%% ─────────────────────────────────────────────────────────────────────
\section{Practical Patterns}\label{sec:enum-practical}

\subsection{Enums and Loops}

Building enum values inside loops and processing them with match is a
common workflow:

\begin{lstlisting}[language=Lattice, caption={Batch-constructing enums}]
enum Status {
    Pending,
    Running(any),
    Done(any),
    Failed(any, any)
}

flux statuses = []
for i in 0..8 {
    let s = match i % 4 {
        0 => Status::Pending,
        1 => Status::Running("task_" + to_string(i)),
        2 => Status::Done(i * 10),
        _ => Status::Failed("err_" + to_string(i), i)
    }
    statuses.push(s)
}

flux pending_count = 0
flux done_count = 0
for s in statuses {
    match s {
        Status::Pending => { pending_count += 1 },
        Status::Done(_) => { done_count += 1 },
        _ => {}
    }
}
print("pending: " + to_string(pending_count))  // pending: 2
print("done: " + to_string(done_count))        // done: 2
\end{lstlisting}

\subsection{Enums as Function Return Types}

Use enums to signal outcomes without exceptions:

\begin{lstlisting}[language=Lattice, caption={Validation with enums}]
enum ValidationResult {
    Valid,
    TooShort(any),
    TooLong(any),
    InvalidChar(any)
}

fn validate_username(name: String) -> ValidationResult {
    if len(name) < 3 {
        return ValidationResult::TooShort(len(name))
    }
    if len(name) > 20 {
        return ValidationResult::TooLong(len(name))
    }
    return ValidationResult::Valid
}

let result = validate_username("Al")
match result {
    ValidationResult::Valid         => print("Username accepted"),
    ValidationResult::TooShort(n)  => print("Too short: " + to_string(n) + " chars"),
    ValidationResult::TooLong(n)   => print("Too long: " + to_string(n) + " chars"),
    ValidationResult::InvalidChar(c) => print("Bad char: " + to_string(c)),
    _ => print("?")
}
// Output: Too short: 2 chars
\end{lstlisting}

\subsection{Under the Hood}
\index{OP\_BUILD\_ENUM}

At runtime, an enum value is represented as a tagged union (see
\filename{include/value.h}):

\begin{itemize}
  \item \textbf{\texttt{enum\_name}}: the enum's type name (e.g., ``Option'')
  \item \textbf{\texttt{variant\_name}}: the variant (e.g., ``Some'')
  \item \textbf{\texttt{payload}}: an array of values (empty for unit
    variants)
  \item \textbf{\texttt{payload\_count}}: the number of payload elements
\end{itemize}

The compiler emits \texttt{OP\_BUILD\_ENUM} with the enum name, variant name,
and payload count as operands. The VM pops payload values from the stack
and assembles the enum value. You can explore the compilation logic in
\filename{src/stackcompiler.c}---look for the \texttt{EXPR\_ENUM\_VARIANT}
case.

%% ─────────────────────────────────────────────────────────────────────
\section{Exercises}\label{sec:enum-exercises}

\begin{enumerate}
  \item \textbf{Traffic light.}
    Define an enum \lstinline|Light| with variants \lstinline|Red|,
    \lstinline|Yellow|, and \lstinline|Green|. Write a function
    \lstinline|next_light(l: Light) -> Light| that cycles through the
    states: Green → Yellow → Red → Green. Run the cycle 6 times and
    print each state.

  \item \textbf{Shape areas.}
    Define an enum \lstinline|Shape| with variants \lstinline|Circle(any)|,
    \lstinline|Rectangle(any, any)|, and \lstinline|Square(any)|. Write a
    function that returns the area. Use 3.14159 for pi.

  \item \textbf{Linked list.}
    Model a linked list with an enum:
    \lstinline|enum List { Cons(any, any), Nil }|. Write functions
    \lstinline|list_sum| (sum all elements) and \lstinline|list_len|
    (count elements). Build a list of \lstinline|[1, 2, 3, 4, 5]| and
    verify both functions.

  \item \textbf{Result pipeline.}
    Write a function \lstinline|map_result(r: Result, f: Fn) -> Result|
    that applies \lstinline|f| to the value inside \lstinline|Result::Ok|
    and passes \lstinline|Result::Err| through unchanged. Chain three
    \lstinline|map_result| calls to build a pipeline.

  \item \textbf{Exhaustiveness experiment.}
    Define an enum with five variants. Write a match expression that covers
    only three of them (no wildcard). Compile the program and observe the
    compiler warning. Then fix it by adding the missing variants.
\end{enumerate}

%% ─────────────────────────────────────────────────────────────────────
\section*{What's Next}

We have seen how to group fields into structs and model alternatives with
enums. But so far, attaching methods to a struct has required declaring a
trait---even when the trait is implemented by only one type. In the next
chapter we explore \emph{traits and impl blocks} in depth: how to declare
interfaces, implement them for multiple types, and use trait-driven
polymorphism to write code that works across any struct that speaks the
right protocol.
